\chapter{DNS interno e intranet de Recursos Humanos}

Responsable: Andrés Rafael Chivalán. Este capítulo describe el servicio de resolución de
nombres interno y la intranet de Recursos Humanos alojados en la subred de Data Center
\code{10.10.0.160/28}. Ambos servicios son privados: no tienen IPv4 pública, no tienen ruta
hacia Internet y solo aceptan tráfico de las redes internas autorizadas.

\section{Componentes}

\begin{table}[H]
\centering
\caption{Servicios de Data Center}
\small
\begin{tabularx}{\textwidth}{>{\raggedright\arraybackslash}p{2.5cm}p{4.3cm}p{2.1cm}>{\raggedright\arraybackslash}p{2.0cm}Y}
\toprule
\textbf{Servicio} & \textbf{Nombre} & \textbf{IP} & \textbf{Puerto} & \textbf{Software} \\
\midrule
DNS interno & \code{ns1.rednova.internal} & \code{10.10.0.164} & UDP/TCP 53 & BIND9, Amazon Linux 2023 \\
Intranet RR. HH. & \code{rrhh.rednova.internal} & \code{10.10.0.165} & TCP 80 & Nginx, HTML estático \\
\bottomrule
\end{tabularx}
\end{table}

Se eligió el dominio \code{rednova.internal} en lugar de \code{.local}, porque este último está
reservado para DNS de multidifusión (mDNS, RFC 6762) y provoca resoluciones inconsistentes en
clientes que implementan ese protocolo.

\section{Funcionamiento del DNS interno}

BIND9 cumple dos funciones. Es \emph{autoritativo} para la zona directa
\code{rednova.internal} y para la zona inversa \code{0.10.10.in-addr.arpa}. Además es
\emph{recursivo} para las redes autorizadas: los nombres que no pertenecen a la zona interna se
reenvían al resolvedor de la VPC (\code{10.10.0.2}, Route 53 Resolver) con
\code{forward only}. Así, un cliente interno puede usar BIND9 sin perder la resolución de
nombres públicos.

\begin{table}[H]
\centering
\caption{Registros de la zona \code{rednova.internal}}
\begin{tabularx}{\textwidth}{p{3.2cm}C{1.6cm}p{4.2cm}Y}
\toprule
\textbf{Nombre} & \textbf{Tipo} & \textbf{Valor} & \textbf{Uso} \\
\midrule
\code{ns1} & A & \code{10.10.0.164} & Servidor DNS autoritativo \\
\code{dns} & CNAME & \code{ns1} & Alias del servidor DNS \\
\code{rrhh} & A & \code{10.10.0.165} & Intranet de Recursos Humanos \\
\code{intranet} & CNAME & \code{rrhh} & Alias de la intranet \\
\code{ti-cliente} & A & \code{10.10.0.132} & Cliente de pruebas de TI \\
\code{web-dmz} & A & \code{10.10.0.196} & Nombre interno de la web pública \\
\bottomrule
\end{tabularx}
\end{table}

\subsection{DNS dividido}

No se modificó el \emph{DHCP Options Set} de la VPC, porque ese cambio afectaría también a
Visitas y a la DMZ, cuyas políticas solo permiten consultar \code{10.10.0.2}. En su lugar, cada
cliente interno recibe un archivo de \code{systemd-resolved} con
\code{DNS=10.10.0.164} y \texttt{Domains=\textasciitilde{}rednova.internal}. Con esa
configuración únicamente los nombres de la zona interna se envían a BIND9 y el resto sigue el
resolvedor entregado por DHCP.

\section{Funcionamiento de la intranet}

Nginx publica una página HTML estática identificada como intranet de Recursos Humanos. El
servidor virtual responde a \code{rrhh.rednova.internal}, \code{intranet.rednova.internal} y
a la IP \code{10.10.0.165}; cualquier otro encabezado \code{Host} recibe el código 444 y la
conexión se cierra. La página no carga fuentes, scripts ni imágenes externas, lo cual es
coherente con una subred sin salida a Internet.

\section{Control de acceso en capas}

\begin{table}[H]
\centering
\caption{Controles aplicados al DNS y a la intranet}
\begin{tabularx}{\textwidth}{p{3.4cm}p{3.4cm}Y}
\toprule
\textbf{Capa} & \textbf{Mecanismo} & \textbf{Regla} \\
\midrule
Red AWS & Security Group \code{dns} & UDP/TCP 53 solo desde Ventas, TI, Data Center y Administración \\
Red AWS & Security Group \code{intranet} & TCP 80 solo desde las mismas cuatro redes \\
Red AWS & Subred privada & Sin IP pública ni ruta \code{0.0.0.0/0} \\
Red AWS & NACL de Visitas y DMZ & Denegación de salida hacia \code{10.10.0.0/24} \\
Servicio & ACL de BIND9 & \code{allow-query} y \code{allow-recursion} limitados a \code{redes_autorizadas}; otros orígenes reciben \code{REFUSED} \\
Servicio & \code{allow}/\code{deny} de Nginx & Mismas cuatro redes; cualquier otro origen recibe 403 \\
Administración & Security Group & SSH únicamente desde el Security Group del bastión \\
\bottomrule
\end{tabularx}
\end{table}

Las reglas del servicio repiten las del Security Group a propósito. Si una regla de AWS se
amplía por error, el servicio sigue rechazando los orígenes no autorizados.

\section{Hardening aplicado}

\begin{itemize}
  \item BIND9 escucha solo en \code{127.0.0.1} y \code{10.10.0.164}; IPv6 deshabilitado.
  \item Transferencias de zona y actualizaciones dinámicas deshabilitadas.
  \item Versión, nombre de host e identificador del servidor ocultos en las respuestas.
  \item Limitación de tasa de respuestas (\code{rate-limit}) para reducir amplificación.
  \item Registro de consultas habilitado como evidencia de auditoría.
  \item Nginx con \code{server_tokens off}, solo métodos GET y HEAD, archivos ocultos denegados
        y cabeceras \code{X-Frame-Options}, \code{X-Content-Type-Options},
        \code{Referrer-Policy} y \code{Content-Security-Policy}.
  \item Archivos de configuración propiedad de \code{root} con permisos \code{0640} o \code{0644}.
  \item SSH solo con llave, sin acceso directo de \code{root} y con tres intentos como máximo.
  \item IMDSv2 obligatorio y volumen EBS cifrado en ambas instancias.
\end{itemize}

\section{Instalación sin NAT Gateway}

Data Center es una subred privada sin NAT, por lo que \code{dnf} no puede descargar paquetes
desde Internet. Los repositorios de Amazon Linux 2023 se alojan en buckets S3 administrados por
AWS. Por ello se agregó un \emph{VPC Gateway Endpoint} de S3 asociado a la tabla de rutas
privada. Este tipo de endpoint no tiene costo y su política solo permite leer el bucket de
repositorios de la región. Los Security Groups del DNS y de la intranet agregan una única
salida TCP 443 hacia la lista de prefijos de S3.

\begin{warningbox}[Alternativa descartada]
Un NAT Gateway habría resuelto la instalación, pero tiene costo por hora y por GB procesado y
abriría salida general a Internet desde redes que no la necesitan.
\end{warningbox}

\section{Automatización con Terraform}

El archivo \code{services.tf} crea ambas instancias con IP privada fija, sin IPv4 pública y
con \code{user_data} generado desde \code{templates/}. Los archivos de \code{services/dns/} y
\code{services/intranet/} se incrustan en base64, de modo que la instancia no depende de Git ni
de Internet. Durante el arranque el script instala los paquetes, copia las configuraciones,
ejecuta \code{named-checkconf}, \code{named-checkzone} y \code{nginx -t} y solo entonces
habilita los servicios. Tres precondiciones detienen el plan si cambian las IP o el dominio sin
actualizar las zonas, o si se activa la Fase 2 sin el bastión de la Fase 1.

\section{Validación local}

El script \code{scripts/validate-andres-phase2.py} comprueba la coherencia entre BIND9, Nginx,
Terraform y el plan de direccionamiento, y ejecuta los validadores nativos disponibles. La
salida se guarda en \code{evidence/fase-2/dns/andres-validation.txt}.

\begin{infobox}[Resultado local]
28 comprobaciones estructurales correctas, 3 omitidas y 0 fallidas. El entorno Windows de
integración no dispone de \code{named-checkconf}, \code{named-checkzone} ni \code{nginx}; por ello esas pruebas
nativas se ejecutan en el \code{user_data} antes de habilitar cada servicio. Las respuestas DNS,
HTTP y los bloqueos se validaron después en AWS mediante la suite de aceptación.
\end{infobox}

\section{Pruebas en AWS}

\begin{table}[H]
\centering
\caption{Pruebas del DNS interno y la intranet}
\small
\begin{tabularx}{\textwidth}{C{1.1cm}p{2.3cm}p{6.4cm}p{2.3cm}Y}
\toprule
\textbf{ID} & \textbf{Origen} & \textbf{Comando} & \textbf{Esperado} & \textbf{Resultado} \\
\midrule
F2-01 & Ventas & \code{dnsquery.py rrhh.rednova.internal} & \code{10.10.0.165} & OK: resuelto \\
F2-02 & Ventas & \code{curl http://rrhh.rednova.internal/} & HTTP 200 & OK: HTTP 200 \\
F2-03 & TI & DNS y HTTP por nombre interno & Resolución y HTTP 200 & OK: 2 de 2 \\
F2-04 & Administración & \code{curl http://10.10.0.165/} & HTTP 200 & OK: HTTP 200 \\
F2-05 & Visitas y DMZ & DNS, HTTP y SSH internos & Filtrado o tiempo agotado & OK: 4 de 4 \\
\bottomrule
\end{tabularx}
\end{table}

La consulta de Ventas recibió una respuesta autoritativa de \code{10.10.0.164} y resolvió
\path{rrhh.rednova.internal} como \code{10.10.0.165}.

\begin{figure}[H]
  \centering
  \includegraphics[width=\textwidth]{05-dns-interno-exitoso.png}
  \caption{Resolución DNS interna exitosa desde una red autorizada. La respuesta contiene un
  registro A para la intranet y el estado \code{NOERROR}.}
  \label{fig:f2-dns-ok}
\end{figure}

\begin{figure}[H]
  \centering
  \includegraphics[width=\textwidth]{06-intranet-interna-exitosa.png}
  \caption{Acceso HTTP exitoso a la intranet de Recursos Humanos desde TI. Nginx devuelve
  \code{HTTP/1.1 200 OK} y el contenido esperado.}
  \label{fig:f2-intranet-ok}
\end{figure}

La prueba negativa desde Visitas agotó el tiempo tanto para el servicio HTTP como para la consulta
al DNS interno. El resultado confirma que la navegación pública permitida a Visitas no abre acceso
lateral hacia Data Center.

\begin{figure}[H]
  \centering
  \includegraphics[width=\textwidth]{07-intranet-bloqueada-desde-visitas.png}
  \caption{Bloqueo desde Visitas hacia la intranet y el DNS interno. El código HTTP \code{000} y
  el tiempo agotado son coherentes con filtrado de red, no con una respuesta del servicio.}
  \label{fig:f2-visitas-bloqueo-interno}
\end{figure}
